\documentclass[../../main.tex]{subfiles}
\begin{document}
\section{流形上的张量与张量场}\label{section:ch12-流形上张量}

现在转到光滑流形。对每个 $p\in M$，切空间 $T_pM$ 是有限维向量空间。逐点地应用上一节构造，可得 $T_pM$ 上的协变、反变与混合张量空间。把它们粘合起来，便得到 $M$ 上的各张量丛。

令 $T^{(k,l)}TM$ 表示 $\binom{k}{l}$ 型张量丛，即
\begin{align*}
T^{(k,l)}TM=\bigsqcup_{p\in M}T^{(k,l)}(T_pM),
\end{align*}
其中
\begin{align*}
T^{(0,k)}TM&=T^kT^*M,\\
T^{(k,0)}TM&=T^kTM.
\end{align*}

\begin{exercise}\label{exercise:ch12-18}
证明 $T^kT^*M$、$T^kTM$ 与 $T^{(k,l)}TM$ 具有 $M$ 上光滑向量丛的自然结构，并确定它们的秩。
\end{exercise}

这些丛中的任一个都称为 $M$ 上的\textbf{张量丛}。（故切丛与余切丛是张量丛的特例。）张量丛的截面称为 $M$ 上的（协变、反变或混合）张量场。光滑张量场即张量丛光滑截面意义下的光滑截面。利用上述等同，我们看到反变 $1$ 张量场就是向量场，协变 $1$ 张量场就是余向量场。由于 $0$ 张量就是实数，$0$ 张量场就是连续实值函数。

这些张量丛的光滑截面空间 $\Gamma(T^kT^*M)$、$\Gamma(T^kTM)$、$\Gamma(T^{(k,l)}TM)$ 是 $\mathbb{R}$ 上的无限维向量空间，也是 $C^\infty(M)$ 上的模。在任意光滑局部坐标 $(x^i)$ 中，这些丛的截面可写为（求和约定生效）：
\begin{align*}
A&=A_{i_1\dots i_k}\,dx^{i_1}\otimes\cdots\otimes dx^{i_k}, && A\in\Gamma(T^kT^*M),\\
A&=A^{i_1\dots i_k}\frac{\partial}{\partial x^{i_1}}\otimes\cdots\otimes\frac{\partial}{\partial x^{i_k}}, && A\in\Gamma(T^kTM),\\
A&=A^{i_1\dots i_k}_{j_1\dots j_l}\frac{\partial}{\partial x^{i_1}}\otimes\cdots\otimes\frac{\partial}{\partial x^{i_k}}\otimes dx^{j_1}\otimes\cdots\otimes dx^{j_l}, && A\in\Gamma(T^{(k,l)}TM).
\end{align*}
函数 $A_{i_1\dots i_k}$、$A^{i_1\dots i_k}$ 或 $A^{i_1\dots i_k}_{j_1\dots j_l}$ 称为 $A$ 在所选坐标中的\textbf{分量函数}。由于我们主要关注光滑协变张量场，对所有光滑协变 $k$ 张量场的空间采用如下简写记号：
\begin{align*}
\mathcal T^k(M)=\Gamma(T^kT^*M).
\end{align*}

\begin{proposition}[张量场的光滑性准则]\label{proposition:ch12-张量场光滑准则}
设 $M$ 为带边或不带边光滑流形，$A\colon M\to T^kT^*M$ 为粗截面。下列条件等价：
\begin{enumerate}[(a)]
\item $A$ 是光滑的；
\item 在每个光滑坐标图中，$A$ 的分量函数光滑；
\item $M$ 的每一点都含于某个坐标图，在其中 $A$ 的分量函数光滑；
\item 若 $X_1,\dots,X_k\in\mathfrak X(M)$，则由 $A(X_1,\dots,X_k)(p)=A_p(X_1|_p,\dots,X_k|_p)$ 定义的函数 $A(X_1,\dots,X_k)\colon M\to\mathbb{R}$ 光滑；
\item 每当 $X_1,\dots,X_k$ 是定义在某开集 $U\subseteq M$ 上的光滑向量场，函数 $A(X_1,\dots,X_k)$ 在 $U$ 上光滑。
\end{enumerate}
\end{proposition}

\begin{exercise}\label{exercise:ch12-20}
证明命题~\ref{proposition:ch12-张量场光滑准则}。
\end{exercise}

\begin{exercise}\label{exercise:ch12-21}
对反变与混合张量场，表述并证明类似于命题~\ref{proposition:ch12-张量场光滑准则}~的光滑性准则。
\end{exercise}

\begin{proposition}\label{proposition:ch12-张量场乘积}
设 $M$ 为带边或不带边光滑流形，$A\in\mathcal T^k(M)$，$B\in\mathcal T^l(M)$，$f\in C^\infty(M)$。则 $fA$ 与 $A\otimes B$ 也是光滑张量场，它们在任意光滑局部坐标图中的分量为
\begin{align*}
(fA)_{i_1\dots i_k}=fA_{i_1\dots i_k},\qquad
(A\otimes B)_{i_1\dots i_{k+l}}=A_{i_1\dots i_k}B_{i_{k+1}\dots i_{k+l}}.
\end{align*}
\end{proposition}

\begin{exercise}\label{exercise:ch12-23}
证明命题~\ref{proposition:ch12-张量场乘积}。
\end{exercise}

命题~\ref{proposition:ch12-张量场光滑准则}(d) 表明：若 $A$ 是 $M$ 上光滑协变 $k$ 张量场，$X_1,\dots,X_k$ 是光滑向量场，则 $A(X_1,\dots,X_k)$ 是 $M$ 上光滑实值函数。故 $A$ 诱导映射
\begin{align*}
\underbrace{\mathfrak X(M)\times\cdots\times\mathfrak X(M)}_{k\text{ 个}}\to C^\infty(M).
\end{align*}
容易看到此映射对 $\mathbb{R}$ 多重线性。事实上更强：它对 $C^\infty(M)$ 多重线性，即对 $f,f'\in C^\infty(M)$ 与 $X_i,X_i'\in\mathfrak X(M)$，有
\begin{align*}
A(X_1,\dots,fX_i+f'X_i',\dots,X_k)=fA(X_1,\dots,X_i,\dots,X_k)+f'A(X_1,\dots,X_i',\dots,X_k).
\end{align*}
下一引理表明此性质是光滑张量场的特征。

\begin{lemma}[张量刻画引理]\label{lemma:ch12-张量刻画}
一个映射
\begin{align}
A\colon \underbrace{\mathfrak X(M)\times\cdots\times\mathfrak X(M)}_{k\text{ 个}}\to C^\infty(M)\label{align:ch12-张量刻画}
\end{align}
可由如上一个光滑协变 $k$ 张量场诱导，当且仅当它对 $C^\infty(M)$ 多重线性。
\end{lemma}

\begin{proof}
我们已经指出：若 $A$ 是光滑协变 $k$ 张量场，则映射 $(X_1,\dots,X_k)\mapsto A(X_1,\dots,X_k)$ 对 $C^\infty(M)$ 多重线性。为证逆命题，沿用丛同态刻画引理（引理~10.29）证明中的论证。

设 $A$ 是\eqref{align:ch12-张量刻画}~中的映射，且对 $C^\infty(M)$ 多重线性。先证 $A$ 局部作用。若 $X_i$ 是在 $p$ 的某邻域 $U$ 上消失的光滑向量场，可选取支集含于 $U$、在 $p$ 处为 $1$ 的鼓包函数 $\varphi$；于是因 $X_i\equiv 0$，有
\begin{align*}
0=A(X_1,\dots,\varphi X_i,\dots,X_k)(p)=\varphi(p)A(X_1,\dots,X_i,\dots,X_k)(p).
\end{align*}
如引理~10.29~证明那样，由此 $A(X_1,\dots,X_k)$ 在 $p$ 处的值只依赖于 $X_1,\dots,X_k$ 在 $p$ 附近的值。

再证 $A$ 确实逐点作用。若 $X_i|_p=0$，则在以 $p$ 为中心的任意坐标图中可写 $X_i=X_i^j\partial/\partial x^j$，其中分量函数 $X_i^j$ 都在 $p$ 处消失。由向量场扩张引理，可在 $M$ 上找整体光滑向量场 $E_j$，使 $E_j$ 在 $p$ 某邻域等于 $\partial/\partial x^j$；类似地，局部定义的函数 $X_i^j$ 可扩张为 $M$ 上整体光滑函数 $f_i^j$，在 $p$ 某邻域与 $X_i^j$ 一致。由 $A$ 对 $C^\infty(M)$ 的多重线性性及 $f_i^j E_j=X_i$ 在 $p$ 附近，得
\begin{align*}
A(X_1,\dots,X_i,\dots,X_k)(p)=A(X_1,\dots,f_i^j E_j,\dots,X_k)(p)=f_i^j(p)A(X_1,\dots,E_j,\dots,X_k)(p)=0.
\end{align*}
由线性性，$A(X_1,\dots,X_k)$ 只依赖于 $X_i$ 在 $p$ 处的值。

现定义粗张量场 $A\colon M\to T^kT^*M$ 为
\begin{align*}
A_p(v_1,\dots,v_k)=A(V_1,\dots,V_k)(p),
\end{align*}
其中 $p\in M$，$v_1,\dots,v_k\in T_pM$，$V_1,\dots,V_k$ 是 $v_1,\dots,v_k$ 到 $M$ 上整体光滑向量场的任意扩张。上述讨论表明这与扩张选取无关，所得张量场由命题~\ref{proposition:ch12-张量场光滑准则}(d)~光滑。
\end{proof}

流形（带边或不带边）上的对称张量场，就是每点取值为对称张量的协变张量场。两个或多个张量场的对称积逐点定义，一如张量积。例如，若 $A,B$ 是光滑余向量场，它们的对称积是光滑 $2$ 张量场 $AB$，由命题~12.15(b)~给出
\begin{align*}
AB=\frac{1}{2}(A\otimes B+B\otimes A).
\end{align*}
交错张量场称为微分形式，我们将从第~14~章起深入研究。

\subsection*{张量场的拉回}

一如余向量场，协变张量场可由光滑映射拉回，得到定义域上的张量场。（此构造只对协变张量有效。）若 $F\colon M\to N$ 是光滑映射，$B$ 是 $N$ 上的协变 $k$ 张量场，则 $B$ 的拉回 $F^*B$ 定义为
\begin{align*}
(F^*B)_p(v_1,\dots,v_k)=B_{F(p)}(dF_p(v_1),\dots,dF_p(v_k)).
\end{align*}
它与余向量场的拉回、以及函数的微分相容。

\subsection*{张量场的 Lie 导数}

设 $V$ 是 $M$ 上的光滑向量场，$\theta_t$ 是它的流。若 $A$ 是 $M$ 上的光滑协变张量场，我们想度量 $A$ 沿 $V$ 的流如何变化。自然做法是把 $A_{\theta_t(p)}$ 用切映射拉回到 $T_pM$，再对 $t$ 求导。对任意 $p\in M$，当 $t$ 充分接近零时，$\theta_t$ 是从 $p$ 某邻域到 $\theta_t(p)$ 某邻域的微分同胚，故 $d(\theta_t)^*_p$ 把 $\theta_t(p)$ 处的张量拉回到 $p$ 处，公式为
\begin{align*}
d(\theta_t)^*_p\big(A_{\theta_t(p)}\big)(v_1,\dots,v_k)
=A_{\theta_t(p)}\big(d(\theta_t)_p(v_1),\dots,d(\theta_t)_p(v_k)\big).
\end{align*}
注意 $d(\theta_t)^*_p(A_{\theta_t(p)})$ 正是拉回张量场 $\theta_t^*A$ 在 $p$ 处的值。

给定 $M$ 上光滑协变张量场 $A$，定义 $A$ 关于 $V$ 的\textbf{Lie 导数}，记为 $\mathcal L_V A$，为
\begin{align}
(\mathcal L_V A)_p=\frac{\mathrm{d}}{\mathrm{d}t}\bigg|_{t=0}(\theta_t^*A)_p
=\lim_{t\to 0}\frac{d(\theta_t)^*_p(A_{\theta_t(p)})-A_p}{t},\label{align:ch12-张量Lie导数定义}
\end{align}
只要导数存在（图~\ref{figure:ch12-张量Lie导数}）。由于被求导表达式对所有 $t$ 都位于 $T^k(T_p^*M)$ 中，$(\mathcal L_V A)_p$ 作为 $T^k(T_p^*M)$ 的元素有意义。

\begin{lemma}\label{lemma:ch12-Lie导数存在}
在上述 $M,V,A$ 下，\eqref{align:ch12-张量Lie导数定义}~中的导数对每个 $p\in M$ 都存在，并把 $\mathcal L_V A$ 定义为 $M$ 上的光滑张量场。
\end{lemma}

\begin{exercise}\label{exercise:ch12-31}
证明上一引理。
\end{exercise}

\begin{figure}[htbp]
\centering
\includegraphics[width=0.6\textwidth]{fig-12-1.png}
\caption{张量场的 Lie 导数}\label{figure:ch12-张量Lie导数}
\end{figure}

\begin{proposition}\label{proposition:ch12-Lie导数性质}
设 $M$ 为光滑流形，$V\in\mathfrak X(M)$。设 $f$ 是 $M$ 上光滑实值函数（视为 $0$ 张量场），$A,B$ 是 $M$ 上光滑协变张量场。
\begin{enumerate}[(a)]
\item $\mathcal L_V f=Vf$；
\item $\mathcal L_V(fA)=(\mathcal L_V f)A+f\,\mathcal L_V A$；
\item $\mathcal L_V(A\otimes B)=(\mathcal L_V A)\otimes B+A\otimes\mathcal L_V B$；
\item 若 $X_1,\dots,X_k$ 是光滑向量场，$A$ 是光滑 $k$ 张量场，则
\begin{align}
\mathcal L_V\big(A(X_1,\dots,X_k)\big)
=(\mathcal L_V A)(X_1,\dots,X_k)+A(\mathcal L_V X_1,\dots,X_k)+\cdots+A(X_1,\dots,\mathcal L_V X_k).\label{align:ch12-Lie导数求值}
\end{align}
\end{enumerate}
\end{proposition}

\begin{proof}
令 $\theta$ 为 $V$ 的流。对实值函数 $f$，可写 $\theta_t^*f(p)=f(\theta_t(p))=f\circ\theta(p)(t)$。于是 $\mathcal L_V f$ 的定义化为复合函数 $f\circ\theta(p)$ 关于 $t$ 的普通导数。由于 $\theta(p)$ 是 $V$ 的积分曲线，由命题~3.24~得
\begin{align*}
(\mathcal L_V f)(p)=\frac{\mathrm{d}}{\mathrm{d}t}\bigg|_{t=0}f\circ\theta(p)=df_p(\theta(p)'(0))=df_p(V_p)=Vf(p).
\end{align*}
这就证明了 (a)。

其他断言可用定理~9.38~中所用技巧证明：在 $V$ 的常点附近，若 $(u^i)$ 是使 $V=\partial/\partial u^1$ 的坐标，则由定义立即得 $\mathcal L_V$ 对张量场的作用只是对其系数关于 $u^1$ 取偏导，于是 (b)--(d) 都由普通乘积法则推出。由连续性，同样关系在 $V$ 的支集上成立；在支集的补集上，$V$ 的流平凡，关系亦成立。
\end{proof}

此命题的一个推论是下面的公式，它用 Lie 括号与函数的普通方向导数表出任意光滑协变张量场的 Lie 导数，使我们不必先确定流即可计算 Lie 导数。

\begin{corollary}\label{corollary:ch12-Lie导数公式}
若 $V$ 是光滑向量场，$A$ 是光滑协变 $k$ 张量场，则对任意光滑向量场 $X_1,\dots,X_k$，
\begin{align}
(\mathcal L_V A)(X_1,\dots,X_k)=V\big(A(X_1,\dots,X_k)\big)
-A([V,X_1],X_2,\dots,X_k)-\cdots-A(X_1,\dots,X_{k-1},[V,X_k]).\label{align:ch12-Lie导数括号}
\end{align}
\end{corollary}

\begin{proof}
只需在\eqref{align:ch12-Lie导数求值}~中解出 $(\mathcal L_V A)(X_1,\dots,X_k)$，并把 $\mathcal L_V f$ 换成 $Vf$、$\mathcal L_V X_i$ 换成 $[V,X_i]$。
\end{proof}

\begin{corollary}\label{corollary:ch12-LVdf}
若 $f\in C^\infty(M)$，则 $\mathcal L_V(df)=d(\mathcal L_V f)$。
\end{corollary}

\begin{proof}
对任意 $X\in\mathfrak X(M)$，用\eqref{align:ch12-Lie导数括号}~计算
\begin{align*}
(\mathcal L_V df)(X)=V(df(X))-df([V,X])=V(Xf)-[V,X]f=V(Xf)-(VXf-XVf)=XVf=d(Vf)(X)=d(\mathcal L_V f)(X).
\end{align*}
\end{proof}

公式\eqref{align:ch12-Lie导数括号}~的一个缺点是：为计算 $\mathcal L_V A$ 在某点 $p\in M$ 对向量 $v_1,\dots,v_k$ 的作用，必须先把它们扩张为 $p$ 附近的向量场。但命题~\ref{proposition:ch12-Lie导数性质}~与推论~\ref{corollary:ch12-LVdf}~引出一个在坐标中计算光滑张量场 Lie 导数的简便方法，避免此问题，因为任意张量场可局部写为函数乘以恰当 $1$ 形式的张量积的线性组合。下例说明此技巧。

\begin{example}[坐标中计算 Lie 导数]\label{example:ch12-坐标Lie导数}
设 $A$ 是任意光滑协变 $2$ 张量场，$V$ 是光滑向量场。我们在光滑局部坐标 $(x^i)$ 中计算 $\mathcal L_V A$。首先注意
\begin{align*}
\mathcal L_V dx^i=d(\mathcal L_V x^i)=d(Vx^i)=dV^i.
\end{align*}
因此
\begin{align*}
\mathcal L_V A&=\mathcal L_V(A_{ij}\,dx^i\otimes dx^j)\\
&=\mathcal L_V(A_{ij})\,dx^i\otimes dx^j+A_{ij}(\mathcal L_V dx^i)\otimes dx^j+A_{ij}\,dx^i\otimes(\mathcal L_V dx^j)\\
&=VA_{ij}\,dx^i\otimes dx^j+A_{ij}\,dV^i\otimes dx^j+A_{ij}\,dx^i\otimes dV^j\\
&=\left(VA_{ij}+A_{kj}\frac{\partial V^k}{\partial x^i}+A_{ik}\frac{\partial V^k}{\partial x^j}\right)dx^i\otimes dx^j.
\end{align*}
\end{example}

回忆向量场 $W$ 关于 $V$ 的 Lie 导数为零，当且仅当 $W$ 在 $V$ 的流下不变（定理~9.42）。事实上协变张量场的 Lie 导数有完全相同的解释。若 $A$ 是 $M$ 上光滑张量场，$\theta$ 是 $M$ 上的流，称 $A$ 在 $\theta$ 下不变，若对每个 $t$，映射 $\theta_t$ 在其有定义处把 $A$ 拉回自身；更精确地说，这意味着对 $\theta$ 定义域中所有 $(t,p)$ 有
\begin{align*}
d(\theta_t)^*_p(A_{\theta_t(p)})=A_p.
\end{align*}
若 $\theta$ 是整体流，这等价于对所有 $t\in\mathbb{R}$ 有 $\theta_t^*A=A$。

\begin{proposition}\label{proposition:ch12-非零t导数}
设 $M$ 为带边或不带边光滑流形，$V\in\mathfrak X(M)$。若 $\partial M\neq\varnothing$，另设 $V$ 切于 $\partial M$。令 $\theta$ 为 $V$ 的流。对任意光滑协变张量场 $A$ 与 $\theta$ 定义域中任意 $(t_0,p)$，
\begin{align}
\frac{\mathrm{d}}{\mathrm{d}t}\bigg|_{t=t_0}(\theta_t^*A)_p=(\theta_{t_0}^*(\mathcal L_V A))_p.\label{align:ch12-非零t导数}
\end{align}
\end{proposition}

\begin{proof}
展开\eqref{align:ch12-非零t导数}~中拉回定义后，需证
\begin{align*}
\frac{\mathrm{d}}{\mathrm{d}t}\bigg|_{t=t_0}d(\theta_t)^*_p(A_{\theta_t(p)})=d(\theta_{t_0})^*_p\big((\mathcal L_V A)_{\theta_{t_0}(p)}\big).
\end{align*}
如命题~9.41~证明那样，作变量替换 $t=s+t_0$，得
\begin{align*}
\frac{\mathrm{d}}{\mathrm{d}t}\bigg|_{t=t_0}d(\theta_t)^*_p(A_{\theta_t(p)})
=\frac{\mathrm{d}}{\mathrm{d}s}\bigg|_{s=0}d(\theta_{s+t_0})^*_p(A_{\theta_{s+t_0}(p)})
=\frac{\mathrm{d}}{\mathrm{d}s}\bigg|_{s=0}d(\theta_{t_0})^*_p\,d(\theta_s)^*_{\theta_{t_0}(p)}(A_{\theta_s(\theta_{t_0}(p))})\\
=d(\theta_{t_0})^*_p\frac{\mathrm{d}}{\mathrm{d}s}\bigg|_{s=0}d(\theta_s)^*_{\theta_{t_0}(p)}(A_{\theta_s(\theta_{t_0}(p))})
=d(\theta_{t_0})^*_p\big((\mathcal L_V A)_{\theta_{t_0}(p)}\big).
\end{align*}
\end{proof}

\begin{theorem}\label{theorem:ch12-Lie导数零即不变}
设 $M$ 为光滑流形，$V\in\mathfrak X(M)$。光滑协变张量场 $A$ 在 $V$ 的流下不变，当且仅当 $\mathcal L_V A=0$。
\end{theorem}

\begin{exercise}\label{exercise:ch12-38}
证明定理~\ref{theorem:ch12-Lie导数零即不变}。
\end{exercise}
\end{document}
